\chapter{Trades d'offre}\label{ch:s2:supply-trades}

Les Treasuries s'échangent moins cher dans les jours qui précèdent une grande adjudication, puis se redressent après, bien que la date et le montant de chaque adjudication soient annoncés à l'avance. Lou, Yan
et Zhang ont documenté le phénomène et chiffré le coût caché pour le Trésor entre 9 et 18 points de base du montant de chaque adjudication. Les dealers doivent absorber l'offre,
et ils sont rémunérés pour cela. De 2010 à septembre 2026, le rendement à dix ans a augmenté en moyenne de 1.98 points de base sur les cinq jours précédant ses 301 adjudications
($t$ = 3.0), et le rendement à trente ans de 2.56 avant ses 234 ($t$ = 3.6). Sur le marché synthétique de ce chapitre, un trade court en duration avant chaque
adjudication et long après rapporte 3.44 points de base de rendement par adjudication ($t$ = 3.1 sur 239), un ratio de Sharpe de 0.70 par an, tandis que le résultat
d'une seule adjudication oscille de plusieurs dizaines de points de base. La construction est \texttt{firm.supplytrade}.

\section{Concessions d'adjudication}

\begin{definition}[Concession d'adjudication]\label{def:s2:supply-trades:concession}
Une \emph{concession d'adjudication}\index{concession d'adjudication} est la dépréciation d'un titre, et de ceux qui lui sont proches, dans les jours qui précèdent la
vente d'une nouvelle émission~: le prix que le marché exige pour absorber une offre connue, qui se récupère en partie une fois l'émission distribuée.
\end{definition}

La prime d'émission d'une seule obligation d'entreprise (livre~2, chapitre~21) est la même idée, pour un seul émetteur. Fleming, Nguyen et Rosenberg ont étudié les positions des dealers américains
sur les Treasuries de 1990 à 2020. L'émission a été le principal moteur des variations hebdomadaires de l'inventaire des dealers. Ces variations n'ont été que partiellement compensées
dans les semaines voisines, et elles n'ont pas été couvertes de façon significative par des contrats à terme. Les dealers ont été rémunérés du risque d'inventaire par une hausse ultérieure des cours. Après la
crise, à mesure que le \omterm{def:s2:swap-spread-and-asset-swap-trades:bscost}{coût de bilan} augmentait, les dealers ont pris des positions plus réduites et les ont transférées plus rapidement. La participation croissante des fonds d'investissement aux adjudications
a réduit cette rémunération.

Le phénomène apparaît dans les données publiques (voir la \cref{fig:s2:supply-trades:path}). On regroupe chaque adjudication de bons et d'obligations du Trésor de 2010 à 2026 avec son indice de référence le plus proche
et l'on mesure le rendement à échéance constante de cet indice autour de la date de l'adjudication. Sur l'ensemble des 1\,338 adjudications, ce rendement est supérieur de 1.47 points de base
le jour de l'adjudication par rapport à dix jours plus tôt, et 5 jours après, il a restitué l'essentiel de cet écart~: 0.40 au-dessus. Le phénomène varie selon l'échéance~:

\begin{center}
\small
\begin{tabular}{@{}lccccc@{}}
\toprule
2010--2026 & adjudications & 5 jours avant (pb) & $t$ & 5 jours après (pb) & $t$\\
\midrule
2 ans & 354 & $-0.40$ & $-0.9$ & $-0.62$ & $-1.3$\\
5 ans & 249 & 1.03 & 1.5 & $-2.16$ & $-3.1$\\
7 ans & 200 & $-0.53$ & $-0.7$ & $-2.02$ & $-2.1$\\
10 ans & 301 & 1.98 & 3.0 & $-0.28$ & $-0.4$\\
30 ans & 234 & 2.56 & 3.6 & $-0.81$ & $-1.1$\\
\bottomrule
\end{tabular}
\end{center}

La partie longue de la courbe se déprécie avant, et les 5 et 7 ans s'apprécient après. Le 2 ans ne montre rien. Classée selon la taille de l'émission, la hausse du 10 ans avant
l'adjudication croît avec la taille~: 1.48 points de base pour des adjudications de 14.5~milliards de dollars en moyenne, 1.69 pour 22.3~milliards de dollars, 2.95 pour 36.2~milliards de dollars. Son
repli après, lui, ne suit pas cette tendance~: $-1.64$, $-0.24$ et $+1.44$. Le rendement à échéance constante est une série de fin de séance, si bien que le jour de l'adjudication elle intègre déjà le
résultat de l'adjudication.

\begin{omfigure}
\begin{tikzpicture}
\begin{axis}[width=10cm,height=5.4cm,xlabel={\small jours de séance depuis l'adjudication},ylabel={\small variation du rendement (pb)},tick label style={font=\footnotesize},
  xmin=-10,xmax=10,ymin=-0.5,ymax=2,grid=major,grid style={gray!25},table/col sep=comma]
\addplot[omThm,thick,mark=*,mark size=1.2pt] table[x=day,y=bp]{figdata/strategies-2/13-supply-trades/path.csv};
\addplot[black,thin] coordinates {(0,-0.5) (0,2)};
\end{axis}
\end{tikzpicture}

\omcaption{Trajectoire moyenne du rendement de l'indice de référence du Trésor autour de 1\,338 adjudications de bons et d'obligations, 2010--2026, par rapport à dix jours de séance avant,
d'après les registres d'adjudications de TreasuryDirect et les rendements à échéance constante publiés par la Fed sur FRED. Statistiques dérivées uniquement. Données~:
\texttt{s2\_fetch\_auctions}.}\label{fig:s2:supply-trades:path}
\end{omfigure}

\section{Primes d'émission et primes on-the-run}

\begin{definition}[Prime on-the-run]\label{def:s2:supply-trades:otr}
La \emph{prime on-the-run}\index{prime on-the-run} est l'écart selon lequel le titre le plus récemment émis à une échéance donnée s'échange plus cher que les titres plus anciens
de maturité voisine, en raison de sa liquidité et de sa valeur en repo~; elle s'amenuise à mesure que l'émission vieillit et qu'une plus récente la remplace.
\end{definition}

La prime détermine ce que détient un \omterm{def:s2:supply-trades:trade}{trade d'offre}. La nouvelle émission est bon marché par rapport à la courbe avant son adjudication et devient le on-the-run, plus cher que le reste de la courbe, après celle-ci. Le
titre qu'elle remplace perd son statut et se déprécie. La courbe ajustée de la Fed (chapitre~10) écarte les titres on-the-run et les premiers off-the-run pour cette
raison. Un roll trade vend l'ancien titre on-the-run et achète le nouveau autour de l'adjudication. Il parie que la prime passe de l'un à l'autre
plus lentement que ne l'intègre le marché.

\section{Émissions d'entreprises}

Les émissions d'entreprises fonctionnent de la même manière, avec des opérations plus petites et moins prévisibles. Les émetteurs annoncent le matin et fixent leur prix dans l'après-midi, et les dealers
et les investisseurs achètent avec une concession sur les obligations existantes de l'émetteur. Les semaines de forte émission déprécient l'ensemble du secteur. Avant une grande opération, les
obligations propres de l'émetteur et leurs substituts proches sont vendus, ou couverts par des swaps et des Treasuries, et les couvertures sont levées ensuite. L'effet est prévisible lorsque le
calendrier l'est. Il souffre aussi d'encombrement (crowding)~: de nombreux desks connaissent ce calendrier.

\begin{definition}[Trade d'offre]\label{def:s2:supply-trades:trade}
Un \emph{trade d'offre}\index{trade d'offre} prend le côté opposé d'une offre connue et programmée~: vendeur (ou en sous-pondération) des titres sur le point d'être émis
et de leurs substituts proches avant la vente, acheteur après celle-ci, rémunéré par le besoin qu'ont l'émetteur et les dealers de distribuer l'offre.
\end{definition}

\section{Trader autour du calendrier}

\texttt{firm.supplytrade} rend la concession explicite (voir le \cref{lst:s2:supply-trades:model}). Un rendement de référence comporte un bruit de 5.5 points de base par jour
et une adjudication mensuelle de 20 à 45~milliards de dollars. La capacité des dealers à absorber chaque adjudication varie (une loi lognormale de dispersion 0.4). Dans les cinq
jours qui précèdent une adjudication, le rendement monte de 0.07 point de base par milliard à capacité moyenne, soit 2.60 points de base en moyenne, et dans les cinq jours qui suivent il restitue
70~\% de la hausse. Le trade est court en duration de cinq jours avant jusqu'à la clôture de l'adjudication, et long pendant cinq jours après, en payant 0.1 point de base de
rendement par jambe sur les contrats à terme.

Sur vingt ans et 239 adjudications, il rapporte 3.44 points de base par adjudication ($t$ = 3.1)~: 1.71 provenant de la jambe avant et 1.72 de la jambe après. À 12
adjudications par an, cela donne un ratio de Sharpe de 0.70, et 0.50 pour chaque jambe prise isolément. Classée par taille, la moyenne passe de 1.66 points de base dans le
tiers le plus petit à 4.94 dans le tiers médian et 3.72 dans le plus grand. Chaque milliard de taille ajoute en moyenne 0.24 point de base au trade
(voir la \cref{fig:s2:supply-trades:scatter}). Le bruit est bien plus grand que l'effet~: les adjudications prises isolément varient de $-47.9$ à $+52.7$ points de base. La capacité
des dealers, que le trader ne voit pas, pèse autant que la taille. La pente en taille rapportée à la capacité vaut 0.16 par unité.

\begin{omfigure}
\begin{tikzpicture}
\begin{axis}[width=10cm,height=5.8cm,xlabel={\small taille de l'adjudication (milliards de \$)},ylabel={\small P\&L du trade (pb de rendement)},tick label style={font=\footnotesize},
  xmin=19,xmax=46,ymin=-55,ymax=60,grid=major,grid style={gray!25},table/col sep=comma]
\addplot[only marks,mark=*,mark size=1pt,omThm,opacity=0.6] table[x=size,y=pnl]{figdata/strategies-2/13-supply-trades/scatter.csv};
\addplot[black,thick] table[x=size,y=fit]{figdata/strategies-2/13-supply-trades/fit.csv};
\end{axis}
\end{tikzpicture}

\omcaption{P\&L synthétique du \omterm{def:s2:supply-trades:trade}{trade d'offre} par adjudication en fonction de la taille de l'adjudication, 239 adjudications sur vingt ans, avec la droite des moindres carrés (0.24
point de base par milliard). Données~: \texttt{s2\_supply.results}.}\label{fig:s2:supply-trades:scatter}
\end{omfigure}

\section{Fiches de stratégie}

\begin{strategyfile}{Achat avant l'adjudication, vente après}\label{strat:s2:supply-trades:auction}
\sfield{Qui vous paie, et pourquoi} Le Trésor, par la concession qu'il verse pour distribuer l'offre~; les teneurs de marché disposant d'une capacité limitée.
\sfield{Instruments et plateformes} Les contrats à terme sur Treasuries ou le titre de référence et ses voisins~; le marché when-issued.
\sfield{Signal} Le calendrier des adjudications et leur taille~; le positionnement des dealers lorsqu'il est observable.
\sfield{Dimensionnement et exécution} De petites positions par adjudication, nombreuses adjudications~; les contrats à terme pour un coût faible.
\sfield{Coûts} Les frais des contrats à terme et le bid--ask~; le roll.
\sfield{Comment la stratégie meurt} Davantage de demande hors dealers lors des adjudications~; l'encombrement dû aux autres acteurs qui négocient le calendrier.
\sfield{Horizon, capacité, infrastructure} Quelques jours~; le calendrier.
\sfield{Backtest honnête} Les rendements aux heures réelles des transactions, et non en fin de journée.
\sfield{Sources} Lou, Yan et Zhang (2013)~; rendements 2010--2026~: $+1.98$ point de base avant les adjudications à 10 ans~; le présent chapitre~: 3.44 par adjudication, ratio de Sharpe
0.70.
\end{strategyfile}

\begin{strategyfile}{Roll on-the-run}\label{strat:s2:supply-trades:roll}
\sfield{Qui vous paie, et pourquoi} Les détenteurs qui paient pour la liquidité on-the-run et la valeur du repo.
\sfield{Instruments et plateformes} Les émissions on-the-run nouvelle et ancienne.
\sfield{Signal} Le spread de rendement entre elles, comparé à son profil au cours du cycle.
\sfield{Dimensionnement et exécution} Position équilibrée en DV01~; la vente à découvert financée en reverse repo.
\sfield{Coûts} Le special de l'émission vendue à découvert.
\sfield{Comment la stratégie meurt} Une fuite vers la liquidité élargit la prime.
\sfield{Horizon, capacité, infrastructure} Un cycle d'adjudications.
\sfield{Backtest honnête} Les taux de repo spéciaux.
\sfield{Sources} Aucun chiffre de performance vérifié.
\end{strategyfile}

\begin{strategyfile}{Prime d'émission des obligations d'entreprise}\label{strat:s2:supply-trades:corporate}
\sfield{Qui vous paie, et pourquoi} Les émetteurs qui fixent le prix des nouvelles obligations en dessous de celui de leurs obligations existantes pour les placer rapidement.
\sfield{Instruments et plateformes} Les nouvelles obligations d'entreprise à l'allocation~; les obligations secondaires de l'émetteur.
\sfield{Signal} La concession comparée à la courbe de l'émetteur.
\sfield{Dimensionnement et exécution} Les allocations~; couvertes par des Treasuries ou des swaps.
\sfield{Coûts} Des allocations inférieures aux ordres~; les spreads secondaires.
\sfield{Comment la stratégie meurt} Des livres encombrés qui réduisent la concession.
\sfield{Horizon, capacité, infrastructure} Quelques jours~; les relations avec les souscripteurs.
\sfield{Backtest honnête} Les allocations réelles, et non les ordres.
\sfield{Sources} Aucun chiffre de performance vérifié.
\end{strategyfile}

\begin{strategyfile}{Extension de fin de mois de l'indice}\label{strat:s2:supply-trades:extension}
\sfield{Qui vous paie, et pourquoi} Les investisseurs suiveurs d'indice qui doivent ajouter de la duration lorsque la duration de l'indice augmente en fin de mois, à mesure que de nouvelles émissions y entrent.
\sfield{Instruments et plateformes} Obligations longues ou contrats à terme avant la fin du mois.
\sfield{Signal} L'extension annoncée de l'indice.
\sfield{Dimensionnement et exécution} Positions prises avant le rééquilibrage, dénouées après.
\sfield{Coûts} Les frais des contrats à terme.
\sfield{Comment la stratégie meurt} Beaucoup de desks négocient le même calendrier.
\sfield{Horizon, capacité, infrastructure} Quelques jours.
\sfield{Backtest honnête} Les chiffres d'extension tels que publiés avant la fin du mois.
\sfield{Sources} Aucun chiffre vérifié.
\end{strategyfile}

\section{Tutoriel~: payer pour absorber}\label{tut:s2:supply-trades:tut}

\begin{tutorial}
\textbf{Objectif.} Simuler un calendrier d'adjudications avec des concessions introduites à dessein, négocier ce calendrier et mesurer la configuration réelle autour des adjudications 2010--2026.
\textbf{État final~:} les tableaux et les deux figures.
\begin{tutsteps}
\item \textbf{Adjudications et transaction}.
\omcode{code/firm/supplytrade/firm_supplytrade.py}{38}{62}{Un calendrier avec des concessions proportionnelles à la taille et à la capacité, et les deux jambes de la transaction.}\label{lst:s2:supply-trades:model}
\item \textbf{Résultats}.
\omcode{code/strategies-2/13-supply-trades/python/s2_supply.py}{35}{51}{Par jambe, par taille, et en fonction de la taille.}\label{lst:s2:supply-trades:results}
\item \textbf{Exécuter} \texttt{s2\_fetch\_auctions.py} une fois, puis \texttt{results()} et \texttt{fig\_supply.py}.
\end{tutsteps}
\textbf{Pour aller plus loin.} Donner au trader un signal bruité de la capacité des dealers et ne négocier que lorsqu'elle est faible~; ajouter une \omterm{def:s2:supply-trades:otr}{prime on-the-run} qui évolue à
chaque adjudication~; utiliser les rendements intrajournaliers à l'heure de l'adjudication.
\end{tutorial}

\section{Construction~: trades d'offre}\label{bld:s2:supply-trades:supplytrade}

\begin{build}
\bfield{Rôle} Un calendrier d'adjudications avec des concessions proportionnelles à la taille rapportée à la capacité des dealers, et la transaction qui l'entoure.
\bfield{Interface} \texttt{SupplyConfig(\dots)}, \texttt{simulate\_auctions(cfg)}, \texttt{supply\_trade(sim, cfg, pre, post)}.
\bfield{Règles} La concession se constitue pendant les cinq jours précédents et est en partie restituée pendant les cinq suivants~; coûts par jambe.
\bfield{Tests d'acceptation} \texttt{code/firm/supplytrade/tests/}~: sans bruit, la transaction rapporte la concession et sa restitution~; coûts par jambe~;
concessions proportionnelles à la taille rapportée à la capacité.
\bfield{Pour aller plus loin} Plusieurs échéances~; réouvertures~; une \omterm{def:s2:supply-trades:otr}{prime on-the-run}.
\end{build}

\begin{omsources}
\item D.~Lou, H.~Yan et J.~Zhang, ``Anticipated and repeated shocks in liquid markets'', \emph{Review of Financial Studies} 26(8), 2013.
\item M.~Fleming, G.~Nguyen et J.~Rosenberg, ``How do Treasury dealers manage their positions?'', Federal Reserve Bank of New York Staff Report 299,
revised 2024.
\item TreasuryDirect auction records; Board of Governors H.15 yields, via FRED.
\end{omsources}

\section{Exercices}

\begin{exercise}[$\star$]\label{exo:s2:supply-trades:1}
Calculer la concession d'une adjudication de 40~milliards de dollars à capacité moyenne, sachant qu'elle vaut 0.07 point de base par milliard. Quelle est sa valeur totale~?
\end{exercise}

\begin{exercise}[$\star$]\label{exo:s2:supply-trades:2}
Un trade rapporte 3.44 points de base par adjudication, avec un écart type d'environ 17 points de base, 12 fois par an. Calculer son ratio de Sharpe annuel.
\end{exercise}

\begin{exercise}[$\star$]\label{exo:s2:supply-trades:3}
Lou, Yan et Zhang situent le coût caché entre 9 et 18 points de base de la taille de l'adjudication. Calculer ce coût pour une adjudication de 40~milliards de dollars.
\end{exercise}

\begin{exercise}[$\star\star$]\label{exo:s2:supply-trades:4}
Montrer pourquoi une offre connue et annoncée peut-elle faire bouger les prix~?
\end{exercise}

\begin{exercise}[$\star\star$]\label{exo:s2:supply-trades:5}
Expliquer pourquoi la rémunération pour absorber les adjudications a pu baisser depuis 2008.
\end{exercise}

\begin{exercise}[$\star\star$]\label{exo:s2:supply-trades:6}
Expliquer pourquoi un rendement de fin de journée sous-estime la concession le jour de l'adjudication.
\end{exercise}

\begin{exercise}[$\star\star\star$]\label{exo:s2:supply-trades:7}
\emph{Programmation.} Relancer avec \texttt{SupplyConfig(giveback=0.0)}. Que se passe-t-il pour chaque jambe, et qu'est-ce que cela indique quant à la jambe qui
rémunère les dealers~?
\end{exercise}

\begin{exercise}[$\star\star\star$]\label{exo:s2:supply-trades:8}
\emph{Trouver la faille.} «~Les trois dernières adjudications nous ont chacune coûté 20 points de base, donc l'effet d'offre a disparu.~»
\end{exercise}

\section{Problème~: se faire payer pour absorber}

\begin{problem}[{Problème du week-end --- trader le calendrier des adjudications}]\label{pb:s2:supply-trades:1}
Le calendrier synthétique du chapitre, les adjudications de 2010 à 2026 et les données publiques.

\medskip\noindent\textbf{Partie I --- Concessions.}
\begin{enumerate}
\item Définir la \omterm{def:s2:supply-trades:concession}{concession d'adjudication} et la \omterm{def:s2:supply-trades:otr}{prime on-the-run}.
\item Qu'ont trouvé Lou, Yan et Zhang, et qu'ont-ils estimé~?
\item Qu'ont trouvé Fleming, Nguyen et Rosenberg au sujet des dealers~?
\item Pourquoi la courbe du Board exclut-elle les émissions les plus récentes~?
\end{enumerate}

\medskip\noindent\textbf{Partie II --- Le parcours réel.}
\begin{enumerate}[resume]
\item Donner le parcours moyen autour des adjudications.
\item Quelles maturités se décotent avant, et lesquelles se valorisent après~?
\item Comment le comportement du 10 ans varie-t-il avec la taille~?
\item Que manque la série de fin de journée~?
\end{enumerate}

\medskip\noindent\textbf{Partie III --- Le trade synthétique.}
\begin{enumerate}[resume]
\item Définir le \omterm{def:s2:supply-trades:trade}{trade d'offre}.
\item Décrire la concession synthétique.
\item Donner la moyenne par adjudication du trade, sa $t$-statistique et son ratio de Sharpe, ainsi que ceux de chaque jambe.
\item Pourquoi le résultat d'une seule adjudication est-il si bruité~?
\end{enumerate}

\medskip\noindent\textbf{Partie IV --- Le verdict.}
\begin{enumerate}[resume]
\item Énoncer le \emph{résultat nommé}~: le rendement par adjudication du trade et sa dépendance à la taille de l'adjudication.
\item Qu'apporte la capacité des dealers~?
\item Comment trader l'offre d'obligations d'entreprises~?
\item Quel est le trade d'extension de fin de mois~?
\item Comment backtester honnêtement le trade d'adjudication~?
\item Quelle fiche de stratégie dépend du repo~?
\item Comment ce chapitre se rattache-t-il au chapitre~10~?
\item En une phrase~: qui paie le trader d'offre~?
\end{enumerate}
\end{problem}

\section{Questions d'entretien}

\begin{interviewq}[$\star$ \iqroles{trader}]\label{iq:s2:supply-trades:1}
Pourquoi les Treasuries se décotent-ils avant une adjudication que tout le monde connaît~?
\end{interviewq}

\begin{interviewq}[$\star\star$ \iqroles{chercheur}]\label{iq:s2:supply-trades:2}
Comment mesureriez-vous la \omterm{def:s2:supply-trades:concession}{concession d'adjudication}, et que contrôleriez-vous~?
\end{interviewq}

\begin{interviewq}[$\star\star$ \iqroles{trader}]\label{iq:s2:supply-trades:3}
Vous êtes dealer et soumissionnez à l'adjudication à 10 ans de demain. Comment couvrez-vous votre position~?
\end{interviewq}

\begin{interviewq}[$\star\star$ \iqroles{risque}]\label{iq:s2:supply-trades:4}
Un desk applique chaque mois le trade d'adjudication. Quel est son risque~?
\end{interviewq}

\begin{interviewq}[$\star\star$ \iqroles{développeur}]\label{iq:s2:supply-trades:5}
Concevez un pipeline de données pour les annonces et les résultats des adjudications.
\end{interviewq}

\begin{interviewq}[$\star\star\star$ \iqroles{chercheur}]\label{iq:s2:supply-trades:6}
Avec un bruit quotidien du rendement de $\sigma$ points de base et une concession $c$ construite sur $k$ jours, combien d'adjudications faut-il pour détecter $c$ avec un $t$ de 2~?
\end{interviewq}
